\chapter{Introduction}

\section{Working Dissertation Frame}

Customer-feedback analytics often needs more than a document-level positive or negative judgement. A single review can mention several service features, product issues, support interactions, or user-experience points, and these aspects can carry different sentiment. This dissertation therefore studies aspect-sentiment classification for customer feedback, with a particular focus on what happens when the target aspect taxonomy changes.

The working dissertation frame is structured candidate-label aspect-sentiment classification under taxonomy shift in customer feedback. In this framing, the model is not asked to invent arbitrary topics. Instead, at inference time it receives a supplied set of candidate aspect labels and predicts which candidate aspects are relevant to the review, together with their sentiment. This makes the open-topic claim operational and measurable while matching a realistic analytics setting where organisations may revise or extend their feedback taxonomy.

\section{Current Empirical Position}

The project currently has a set of local and language-model baselines over FABSA. Closed-topic and held-out organisation evaluations establish the fixed-taxonomy reference points. Fixed held-out-aspect and leave-one-aspect-out evaluations then test taxonomy-shift behaviour. The strongest current fixed held-out-aspect local non-LLM baseline is a candidate-aspect DistilBERT selector combined with a DistilBERT aspect-conditioned sentiment classifier, achieving 0.6071 pair samples F1. Under full all-row LOAO, however, the preferred DistilBERT branch drops to 0.3128 mean pair micro F1, showing that the fixed held-out-aspect result overstates robustness across aspect rotations. The Qwen indexed zero-shot LOAO baseline reaches 0.3378 mean test pair micro F1 with stable JSON, but its positive-gold diagnostic is much stronger than its all-row score, indicating that absence calibration on empty-gold rows is the main failure mode. Hosted Gemini and local-to-Gemini cascade experiments are complete for the fixed held-out-aspect setting and provide cost, latency, schema-reliability, and selective-deployment evidence, but they should not be treated as full LOAO robustness results.

\section{Contribution Shape}

The intended contribution is an evaluation-centred study rather than a claim that any one model family is automatically superior. The dissertation compares increasingly difficult generalisation settings: fixed taxonomy, unseen organisations, held-out aspects, and leave-one-aspect-out rotations. It also compares model families that can consume candidate labels, including local discriminative baselines, open-weight structured-output Qwen, hosted Gemini, and selective local-hosted cascades. Full fine-tuned Qwen LoRA LOAO remains the main pending compute-bound experiment; if completed, it will test whether task-specific open-weight adaptation can reduce the zero-shot absence-calibration failure under the same LOAO protocol.
